\scanpage{13}
\mspage{5}

qui avec l'origine $e$ forment l'ens des pts fixes de $X_S$. (L'ens des pts fixes sous $\mathbb Z/6\mathbb Z$ est réduit à $e$.) Les deux pts fixes sous $\mathbb Z/3\mathbb Z$, distincts de $e$, correspondent aux deux triangles \unsure{co-inscrits|coinscrits} de l'hexagone $S$. (\inserted{les deux} triplets des côtés mutuellement non adjacents.)

\uline{Décompositions cellulaires canoniques.} Décomposition en deux faces triangulaires, par trois arêtes, un seul sommet. Les faces triangulaires correspondent aux éléments de $(X_S)^\zeta\setminus\{e\}$ (chaque face en contient un et un seul) i.e. aux triangles \unsure{co-inscrits|coinscrits} de l'hexagone. Les côtés correspondent aux trois

\margin{12}
